\documentclass[11pt]{report}
% define the title
\usepackage[brazil]{babel}
\usepackage[utf8]{inputenc}		%Usar acentos diretos do teclado
\usepackage[T1]{fontenc}		%Padrão para imagens
\usepackage{graphicx}
\usepackage{subfigure}
\usepackage{rotating}
\usepackage{longtable}
\usepackage{amsmath}
\usepackage{amsfonts}
\usepackage{amssymb}
\usepackage{geometry}
\usepackage{caption}
\geometry{left=3cm, right=2cm, top=2cm, bottom=2cm}
% Title Page
\title{Exercícios NLAA}
\author{Pedro Yochinori Gushiken}



\begin{document}
\maketitle
\chapter*{7.1}

\section*{7.1.i}

Seja:
\[
H = I - 2\frac{uu^T}{u^Tu}
\]

Então se

\[
Hu = -u = u - 2\frac{uu^Tu}{u^Tu}
\]

Que leva à:

\[
-2*u^Tuu = -2uu^Tu
\]

Como $u^Tu$ é um escalar a identidade se confirma.

\section*{7.1.ii}

Temos que:

\[
Hv = v
\]

se $v^Tu = 0$. Pre-multiplicando por $u^T$, transpondo, e fazendo uso da identidade anterior temos:

\[
(u^THv)^T = (u^Tv)^T
\]

\[
v^T H u = v^T u
\]

\[
-v^Tu = v^T u
\]

\[
v^Tu = 0
\]

Que confirma a condição imposta para que $Hv = v$ seja verdadeiro.
\chapter*{7.3}
\section*{7.3.i}
Seja uma matriz $A$ de tamanho compatível e uma matriz de Householder conhecida (que implica em u conhecido):

\[
H = I - 2\frac{uu^T}{u^Tu}
\]

O produto $HA$ pode ser realizado da seguinte forma (implicita):
\begin{itemize}
\item[i)] Calcule $v=A^Tu$
\item[ii)]Calcule $\beta = \frac{2}{u^Tu}$
\item[iii)] $HA = A-\beta u v^T$ 
\end{itemize}

Ou da seguinte forma (explicita):

\[
HA = (I - 2\frac{uu^T}{u^Tu}) A
\]

A forma implícita resulta no calculo de um produto de matriz por vetor, seguida de um produto vetor por vetor e finalmente um produto vetor por vetor por escalar em termos de custo computacional.\\

A forma explícita resulta no calculo de uma matriz por outra matriz, tal que o custo computacional da segunda forma é maior.

\section*{7.3.ii}

Similarmente, o cálculo de $AH$ pela forma implícita é:

\begin{itemize}
\item[i)] Calcule $w = Au$
\item[ii)] Calcule $\beta = \frac{2}{u^Tu}$
\item[iii)] $AH = A-\beta wu^T$
\end{itemize}

\section*{7.3.iii}

Desejamos calcular: 

\[
(\mathbf{\prod_{i=1}^{r}}H_i) C
\]

Façamos $C_1 = C$:

\begin{itemize} 
\item[i)] Calcule $\beta_j = \frac{2}{u_j^Tu_j}$
\item[ii)] Calcule $v_j = C_j^T u_j$
\item[iii)] Façamos $C_{j+1} = C_j - \beta_j u_j v_j^T$, de $j=1$ até $j=r$
\item[iv)] Colhemos $(\mathbf{\prod_{i=1}^{r}}H_i) C = C_r$
\end{itemize}

\section*{7.3.iv}

Desejamos calcular o produto do item anterior de forma explicita:

\begin{itemize}
\item[i)] Calcule e armazene as matrizes $H_j = I - 2\frac{u_ju_j^T}{u_j^T u_j}$, de $j=1$ até $j=r$
\item[ii)] Realize o produto $H_1\times H_2\times...\times H_r\times C$
\end{itemize}


\chapter*{9.17}
\section*{Householder}
Seja:

\[
A = \begin{bmatrix}
0 & 1 & 0 & 0 & 0 \\
\mathbf{0} & 0 & 1 & 0 & 0 \\
\mathbf{0} & 0 & 0 & 1 & 0 \\
\mathbf{0} & 0 & 0 & 0 & 1 \\
\mathbf{1} & 2 & 3 & 4 & 5 
\end{bmatrix}
\]

Usando o processo de Householder para obter uma matriz Hessenberg temos:

\[
\hat{P}_1 \begin{bmatrix}
0 \\ 0 \\ 0 \\ 1
\end{bmatrix}
=
\begin{bmatrix}
* \\ 0 \\ 0 \\ 0
\end{bmatrix}
\]

Usando a regra para a escolha de $u$ temos:

\[
u_1 = \begin{bmatrix}
0 \\ 0 \\ 0 \\ 1
\end{bmatrix}
+ ||\begin{bmatrix}
0\\ 0 \\ 0 \\ 1
\end{bmatrix}||_2 e_1
=
\begin{bmatrix}
1\\ 0 \\ 0 \\ 1
\end{bmatrix}
\]

E construímos $P_1$ como:

\[
\hat{P}_1 = I - 2\frac{u_1u_1^T}{u_1^Tu_1}
\]

\[
P_1 = \begin{bmatrix}
I_1 & 0 \\
0   & \hat{P}_1
\end{bmatrix}
=
\begin{bmatrix}
1& 0 & 0 & 0& 0\\
0& 0 & 0 & 0& -1\\
0& 0 & 1 & 0& 0\\
0& 0 & 0 & 1& 0\\
0& -1 & 0 & 0& 0
\end{bmatrix}
\]

E calculamos:

\[
A^{(1)} = P_1AP_1^T = \begin{bmatrix}
0& 0 & 0 & 0& -1\\
-1& -5 & -3 & -4& 2\\
0& \mathbf{0} & 0 & 1& 0\\
0& \mathbf{-1} & 0 & 0& 0\\
0& \mathbf{0} & -1 & 0& 0
\end{bmatrix}
\]

Tal que obtemos o vetor $u_2$:

\[
u = \begin{bmatrix}
0 \\ -1 \\ 0
\end{bmatrix}
+  ||\begin{bmatrix}
0 \\ -1 \\ 0
\end{bmatrix}||_2 e_1 = \begin{bmatrix}
1 \\ -1 \\ 0
\end{bmatrix}
\]

E calculamos:

\[
\hat{P_2} = I_3 - 2\frac{u_2u_2^T}{u_2^Tu_2} = \begin{bmatrix}
0 & 1 & 0 \\
1 & 0 & 0\\
0 & 0 & 1
\end{bmatrix}
\]

Tal que:

\[
P_2 = \begin{bmatrix}
I_2 & 0 \\
0   & \hat{P}_2
\end{bmatrix}
=
\begin{bmatrix}
1& 0 & 0 & 0& 0\\
0& 1 & 0 & 0& 0\\
0& 0 & 0 & 1& 0\\
0& 0 & 1 & 0& 0\\
0& 0 & 0 & 0& 1
\end{bmatrix}
\]

E finalmente obtemos $A_h = P_2 A^{(1)} P_2^T$:

\[
A_h = \begin{bmatrix}
0& 0 & 0 & 0& -1\\
-1& 5 & -4 & -3& 2\\
0& -1 & 0 & 1& 0\\
0& 0 & 1 & 0& 0\\
0& 0 & 0 & -1& 0
\end{bmatrix}
\]

\section*{Givens}

Seja:

\[
A = \begin{bmatrix}
0 & 1 & 0 & 0 & 0 \\
0 & 0 & 1 & 0 & 0 \\
0 & 0 & 0 & 1 & 0 \\
\mathbf{0} & 0 & 0 & 0 & 1 \\
\mathbf{1} & 2 & 3 & 4 & 5 
\end{bmatrix}
\]

Queremos construir uma matriz com norma 1 tal que:

\[
\begin{bmatrix}
c & s \\
-s & c
\end{bmatrix}
\begin{bmatrix}
0 \\ 1
\end{bmatrix}
=
\begin{bmatrix}
* \\
0
\end{bmatrix}
\]

E assim por diante até que a estrutura da primeira coluna seja da primeira coluna de uma matriz de Householder superior, construindo sucessivamente as matrizes de Givens para a primeira série de rotações obtemos:

\[
J_1 = \begin{bmatrix}
1 & 0 & 0 & 0 & 0 \\
0 & 0 & 1 & 0 & 0 \\
0 & -1 & 0 & 0 & 0 \\
0 & 0 & 0 & 1 & 0 \\
0 & 0 & 0 & 0 & 1 
\end{bmatrix}
\]
\[
J_2 = \begin{bmatrix}
1 & 0 & 0 & 0 & 0 \\
0 & 0 & 0 & 1 & 0 \\
0 & 0 & 1 & 0 & 0 \\
0 & -1 & 0 & 0 & 0 \\
0 & 0 & 0 & 0 & 1 
\end{bmatrix}
\]\\
\[
J_3= \begin{bmatrix}
1 & 0 & 0 & 0 & 0 \\
0 & 0 & 0 & 0 & 1 \\
0 & 0 & 1 & 0 & 0 \\
0 & 0 & 0 & 1 & 0 \\
0 &-1 & 0 & 0 & 0 
\end{bmatrix}
\]

Obtemos $A^{(1)} = J_1J_2J_3AJ_1^TJ_2^TJ_3^T$, similarmente ao primeiro passo. tal que:

\[
J_4= \begin{bmatrix}
1 & 0 & 0 & 0 & 0 \\
0 & 1 & 0 & 0 & 0 \\
0 & 0 & 0 & 1 & 0 \\
0 & 0 & -1 & 0 & 0 \\
0 & 0 & 0 & 0 & 1 
\end{bmatrix}
\]

\[
J_5= \begin{bmatrix}
1 & 0 & 0 & 0 & 0 \\
0 & 1 & 0 & 0 & 0 \\
0 & 0 & 0 & 0 & 1 \\
0 & 0 & 0 & 1 & 0 \\
0 & 0 & -1 & 0 & 0 
\end{bmatrix}
\]


Obtemos $A^{(2)} = J_4J_5A^{(1)}J_5^TJ_4^T$, e finalmente:

\[
J_6= \begin{bmatrix}
1 & 0 & 0 & 0 & 0 \\
0 & 1 & 0 & 0 & 0 \\
0 & 0 & 1 & 0 & 0 \\
0 & 0 & 0 & 0 & 1 \\
0 & 0 & 0 & -1 & 0 
\end{bmatrix}
\]

E finalmente, temos $A_h = J_6A^{(2)}J_6^T$

Ambos os métodos produzem matrizes idênticas no sentido de terem autovalores idênticos

\chapter*{11.13}
Seja o par $\{A1,B1\}$:

\[
A_1 = \begin{bmatrix}
1 & 1 & 1\\
1 & 1 & 1\\
1 & 1 & 1
\end{bmatrix}
\]

\[
B_1 = \begin{bmatrix}
10 & 1 & 0\\
1 & 10 & 1\\
0 & 1 & 10
\end{bmatrix}
\]

E o par $\{A2,B2\}$:

\[
A_2 = \begin{bmatrix}
10 & 1 & 1\\
1 & 10 & 1\\
1 & 1 & 1
\end{bmatrix}
\]

\[
B_2 = \begin{bmatrix}
1 & 1/2 & 1/3\\
1/2 & 1/3 & 1/4\\
1/3 & 1/4 & 1/5
\end{bmatrix}
\]

Desejamos encontrar os autovalores e autovetores dos pares através dos métodos pedidos, para tal recorremos ao computador e software MATLAB devido a quantidade de operações que são necessárias:

\begin{verbatim}
%Questão 11.13
A1 = [1 1 1; 1 1 1 ; 1 1 1]; B1 = [10 1 0; 1 10 1; 0 1 10];
A2 = [10 1 1; 1 10 1; 1 1 1]; B2 = [1 1/2 1/3; 1/2 1/3 1/4; 1/3 1/4 1/5];

%a1)
A1 = [1 1 1; 1 1 1 ; 1 1 1]; B1 = [10 1 0; 1 10 1; 0 1 10];
u_1 = [10; 1; 0] + sqrt(101)*[1; 0; 0];
H1  = eye(3)-2*(u_1*u_1')/(u_1'*u_1);
BL = H1*B1;
AL = H1*A1;
u_2 = [9.8509; 1] + sqrt(9.8509^2 +1)*[1; 0];
H2_hat = eye(2) -2*(u_2*u_2')/(u_2'*u_2);
H2 = [1 0 0; 0 H2_hat(1,:); 0 H2_hat(2,:)];
BL = H2*BL;
AL = H2*AL;
s1 = AL(3,1)/AL(2,1);
c = sqrt((s1^2+1)^-1);
s = s1*c;
J1 = [1 0 0; 0 c s; 0 -s c];
AL = J1*AL;
BL = J1*BL;
s1 = BL(3,2)/BL(3,3);
c = sqrt((s1^2+1)^-1);
s = -s1*c;
J2 = [1 0 0; 0 c s; 0 -s c]';
AL1 = AL*J2;
BL1 = BL*J2;
%Fim do primeiro estágio, AL em forma Hessenberg e BL triangular superior

%obs: há um possível erro de digitação na página 458, acredito que a_1,a_2 
%relativos ao cálculo de x e y sejam alfa_1,alfa_2, até mesmo por questão de dimensão
for k=1:4
C = AL1*BL1^-1;
C1 = [C(2,2) C(2,3); C(3,2) C(3,3)];
C_eig = eig(C1);
%Assinalamos os coeficientes para double shift
alfa_1 = C_eig(1,1);
alfa_2 = C_eig(2,1);

C_col12 = [AL1(:,1) AL1(:,2)]*[BL1(1,1) BL1(1,2); 0 BL1(2,2)]^-1;
%Calculamos a primeira coluna de N
x = (C_col12(1,1) - alfa_1)*(C_col12(1,1)-alfa_2) +C_col12(1,2)*C_col12(2,1);
y = C_col12(2,1)*(C_col12(1,1)-alfa_2) + C_col12(2,1)*(C_col12(2,2) - alfa_1);
z = C_col12(2,1)*C_col12(3,2);


n1 = [x; y; z];
u1 = n1 + norm(n1)*[1; 0; 0];
%Encontramos uma matriz de householder apropriada e executamos a transformação
H  = eye(3) - 2*(u1*u1')/(u1'*u1);
AL1 = H*AL1;
BL1 = H*BL1;
%retiramos o termo não nulo indesejado introduzido em B através de um rotação
s1 = BL1(2,1)/BL1(2,2);
c = sqrt((s1^2+1)^-1);
s = -s1*c;
J2 = [c s 0 ; -s c 0; 0 0 1]';
AL1 = AL1*J2;
BL1 = BL1*J2;
end

%tal que os autovalores do feixe são: l1 = 3/11.3, l2 e l3=0
%empregando 11.2 temos:
v0 = [1; 1; 1]; l0 = 3/11.3;
v_hat(:,1) = v0;
for k=1:3
   v_hat(:,k+1) = (AL1-BL1*3/11.3)\BL1*v_hat(:,k);
   v(:,k+1)= v_hat(:,k+1)/norm(v_hat(:,k+1));
end
%que leva a v=[1; 0; 0] associado ao autovalor l1 = 3/11.3
%para os autovalores 0 arbitramos autovetores associados canônicos

A2 = [10 1 1; 1 10 1; 1 1 1]; B2 = [1 1/2 1/3; 1/2 1/3 1/4; 1/3 1/4 1/5];
u_1 = [1; 1/2; 1/3] + sqrt(1+1/4+1/9)*[1; 0; 0];
%matriz de householder 1
H1  = eye(3)-2*(u_1*u_1')/(u_1'*u_1);
BL = H1*B2;
AL = H1*A2;
%matriz de householder 2
u_2 = [0.0696; 0.0742] + sqrt(0.0696^2 + 0.0742^2)*[1; 0];
H2_hat = eye(2) -2*(u_2*u_2')/(u_2'*u_2);
H2 = [1 0 0; 0 H2_hat(1,:); 0 H2_hat(2,:)];
BL = H2*BL;
AL = H2*AL;
%rotação tal que A fique na forma Hessenberg
s1 = AL(3,1)/AL(2,1);
c = sqrt((s1^2+1)^-1);
s = s1*c;
J1 = [1 0 0; 0 c s; 0 -s c];
AL = J1*AL;
BL = J1*BL;
%rotação de forma a restaurar forma triangular superior para B
s1 = BL(3,2)/BL(3,3);
c = sqrt((s1^2+1)^-1);
s = -s1*c;
J2 = [1 0 0; 0 c s; 0 -s c]';
AL2 = AL*J2;
BL2 = BL*J2;
%Fim do primeiro estágio, A em forma Hessenberg superior e B em forma
%triangular superior

for k=1:32
C = AL2*BL2^-1;
C1 = [C(2,2) C(2,3); C(3,2) C(3,3)];
C_eig = eig(C1);
%Computamos os escalares necessários para o double shift
alfa_1 = C_eig(1,1);
alfa_2 = C_eig(2,1);

C_col12 = [AL2(:,1) AL2(:,2)]*[BL2(1,1) BL2(1,2); 0 BL2(2,2)]^-1;
%x,y,z são os elementos não nulos da coluna de N = (C-alfa_1 I)(C-alfa_2 I)
x = (C_col12(1,1) - alfa_1)*(C_col12(1,1)-alfa_2) +C_col12(1,2)*C_col12(2,1);
y = C_col12(2,1)*(C_col12(1,1)-alfa_2) + C_col12(2,1)*(C_col12(2,2) - alfa_1);
z = C_col12(2,1)*C_col12(3,2);


n1 = [x; y; z];
%geramos a matriz de householder apropriada
u1 = n1 + norm(n1)*[1; 0; 0];
H  = eye(3) - 2*(u1*u1')/(u1'*u1);
%transformamos
AL2 = H*AL2;
BL2 = H*BL2;

%restabelecemos as propriedades desejadas
s1 = AL2(2,1)/AL2(3,1);
c = sqrt((s1^2+1)^-1);
s = -s1*c;
J2 = [1 0 0 ; 0 c s; 0 -s c]';
AL2 = J2*AL2;
BL2 = J2*BL2;

s1 = BL2(3,1)/BL2(3,2);
c = sqrt((s1^2+1)^-1);
s = -s1*c;
J2 = [c s 0 ; -s c 0; 0 0 1]';
AL2 = AL2*J2;
BL2 = BL2*J2;
s1 = BL2(3,2)/BL2(3,3);
c = sqrt((s1^2+1)^-1);
s = -s1*c;
J2 = [1 0 0 ; 0 c s; 0 -s c]';
AL2 = AL2*J2;
BL2 = BL2*J2;
s1 = BL2(2,1)/BL2(2,2);
c = sqrt((s1^2+1)^-1);
s = -s1*c;
J2 = [c s 0 ; -s c 0; 0 0 1]';
AL2 = AL2*J2;
BL2 = BL2*J2;
end

%iterando, o vetor converge desde que tenhamos uma
%estimativa apropriada do lambda associado
v0 = [1; 1; 1]; l0 = -0.6109/2.8455;
v_hat(:,1) = v0;
for k=1:80
   v_hat(:,k+1) = (AL1+BL1*0.6109/2.8455)\BL1*v_hat(:,k);
   v(:,k+1)= v_hat(:,k+1)/norm(v_hat(:,k+1));
end

% que leva a [0.015 0.7070 0.7070]^T para o autovalor -0.6109/2.8455


A1 = [1 1 1; 1 1 1 ; 1 1 1]; B1 = [10 1 0; 1 10 1; 0 1 10];
A2 = [10 1 1; 1 10 1; 1 1 1]; B2 = [1 1/2 1/3; 1/2 1/3 1/4; 1/3 1/4 1/5];
%a3) QR Simultâneo usando coeficiente de Wilkinson

%fatorização de cholesky para produzir uma matriz tridiagonal
R1 = chol(B1);

%uso de simetria
C = R1'^-1*A1*(R1)^-1;
T(:,:,1) = C;
for k=1:4
%fatorização QR simultânea, empregando diretamente o algoritmo
r = T(3,3,k)/2 - T(2,2,k)/2;
mu(k) = T(3,3,k) +r -sign(r)*sqrt(r^2 + T(3,2,k)^2);
QR = T(:,:,k) - mu(k)*eye(3);
[Q,R] = qr(QR);
T(:,:,k+1) = R*Q+mu(k)*eye(3);
end

%para o segundo par, mesmo processo
R2 = chol(B2);
C = R2'^-1*A2*(R2)^-1;
T2(:,:,1) = C;
for k=1:800
r = T2(3,3,k)/2 - T2(2,2,k)/2;
mu(k) = T2(3,3,k) +r -sign(r)*sqrt(r^2 + T2(3,2,k)^2);
QR = T2(:,:,k) - mu(k)*eye(3);
[Q,R] = qr(QR);
T2(:,:,k+1) = R*Q+mu(k)*eye(3);
end

%a2) Rayleigh Quotient generalizado iterativo
% o par já é simétrico

A1 = [1 1 1; 1 1 1 ; 1 1 1]; B1 = [10 1 0; 1 10 1; 0 1 10];
A2 = [10 1 1; 1 10 1; 1 1 1]; B2 = [1 1/2 1/3; 1/2 1/3 1/4; 1/3 1/4 1/5];

%escolhemos x0 = [1; 0; 0], que converge para lambda = 0 e falha por erro numérico
x0(:,1) = [1; 0; 0];
for k=1:3
   lambda = x0(:,k)'*A1*x0(:,k)/(x0(:,k)'*B1*x0(:,k)) ;
   x0(:,k+1) = (A1-lambda*B1)\(B1*x0(:,k));
   x0(:,k+1) = x0(:,k+1)/norm(x0(:,k+1));
end

%vários exemplos de x0 para o segundo par
x01(:,1) = [1; 0; 0];
x01(:,1) = [-0.4 ; 0; 0.9];
x01(:,1) = [0.1524 ; 0.6989; -0.6989];
for k=1:4
   lambda1 = x01(:,k)'*A2*x01(:,k)/(x01(:,k)'*B2*x01(:,k)) ;
   x01(:,k+1) = (A2-lambda1*B2)\(B2*x01(:,k));
   x01(:,k+1) = x01(:,k+1)/norm(x01(:,k+1));
end

%nota: este método não é capaz de alcançar o autovalor com valor muito alto, 
%mas escolhas apropriadas retornam o autovalor aproximadamente igual a 4 e 
%o autovalor aproximadamente igual a 20 pois os outros dois autovalores 
%são numericamente atrativos

\end{verbatim}


\chapter*{11.21}

Partimos para o software, usando propriedades para a resolução onde se mostrou possível:

\begin{verbatim}
M1 = eye(3); K1 = [2 -1 0; -1 2 -1; 0 -1 2]; D1=ones(3);
%M1 não singular
A1 = [zeros(3) eye(3); -M1^-1*K1 -M1^-1*D1];
[v1,d1] = eig(A1);
M2 = ones(3); K2 = [10 1 1; 1 10 1; 1 1 10]; D2 = 2*M2+3*K2;
%M2 singular
[v2,d2] = polyeig(M2,K2,D2);
M3 = eye(3); K3 = ones(3); D3=K3;
%M3 não singular
A3 = [zeros(3) eye(3); -M3^-1*K3 -M3^-1*D3];
[v3,d3] = eig(A3);
M4 = [1 2 3; 4 5 6; 7 8 9]; K4=ones(3); D4=zeros(3);
[v4,d4] = polyeig(M4,K4,D4);
M5 = eye(3); K5 = ones(3); D5=zeros(3);
%M5 não singular
A5 = [zeros(3) eye(3); -M5^-1*K5 -M5^-1*D5];
[v5,d5] = eig(A5);
\end{verbatim}

\end{document}         